Abstract

This bachelor’s thesis is devoted to adaptive fault-tolerant data storage in an append-only distributed system. The goal of the work is to develop a model and an architectural approach in which the storage scheme can be changed in the background during the data lifetime instead of being fixed once when an object is written. The thesis analyzes approaches based on replication, erasure coding, and temperature-aware policies; proposes a hybrid storage model with replicas, local parity blocks, and global parity blocks; defines admissible transitions between schemes; and introduces a transition utility model that takes into account data activity, occupied storage space, expected read cost, and the cost of background recoding. An experimental system with a background recoder and planners is also developed. Experiments with data cooling, skewed read workloads, and node failures show that adaptive recoding can use disk space more efficiently than replication while preserving cheaper reads for more active or riskier parts of the stored data compared with using erasure coding permanently.

Keywords: distributed storage; erasure coding; replication; data temperature; background recoding.
